\documentclass[12pt]{article}

% ==================================================
% PAQUETES
% ==================================================
\usepackage[utf8]{inputenc}
\usepackage[T1]{fontenc}
\usepackage[spanish]{babel}
\usepackage{lmodern}
\usepackage{graphicx}
\usepackage{geometry}
\usepackage{hyperref}
\usepackage{float}
\usepackage{array}
\usepackage{booktabs}

% ==================================================
% MÁRGENES
% ==================================================
\geometry{
    a4paper,
    top=2.2cm,
    bottom=2.2cm,
    left=2.5cm,
    right=2.5cm
}

% ==================================================
% HIPERVÍNCULOS
% ==================================================
\hypersetup{
    colorlinks=true,
    linkcolor=blue,
    urlcolor=blue
}

\begin{document}

% ==================================================
% CARÁTULA
% ==================================================

\begin{titlepage}
\centering

% Logo
\includegraphics[width=1.8cm]{logo.jpg}\\[0.4cm]

{\fontsize{14}{17}\selectfont\bfseries
UNIVERSIDAD TÉCNICA ESTATAL DE QUEVEDO
}\\[0.2cm]

{\fontsize{12}{15}\selectfont\bfseries
FACULTAD DE CIENCIAS DE LA INGENIERÍA
}\\[0.15cm]

{\fontsize{12}{15}\selectfont
CARRERA DE INGENIERÍA DE SOFTWARE
}\\[0.6cm]

\rule{0.85\textwidth}{0.5pt}\\[0.4cm]

{\fontsize{17}{20}\selectfont\bfseries
PRUEBA PRÁCTICA -- UNIDAD IV
}\\[0.25cm]

{\fontsize{12}{15}\selectfont\bfseries
VALIDACIÓN, GESTIÓN, HERRAMIENTAS Y ESTÁNDARES DE REQUISITOS
}\\[0.25cm]

{\fontsize{11}{14}\selectfont
\textbf{CASO:} SISTEMA DE GESTIÓN DE PEDIDOS
}\\[0.4cm]

\rule{0.85\textwidth}{0.5pt}\\[0.5cm]

{\fontsize{10.5}{13}\selectfont
\textbf{MATERIA:} Ingeniería de Requisitos\\[0.15cm]

\textbf{CURSO:} 4to A -- Software\\[0.15cm]

\textbf{ESTUDIANTE:} HERRERA RAMOS THAIS MELANIE\\[0.15cm]

\textbf{DOCENTE:} Ing. Gleiston Guerrero, Mg.
}\\[0.5cm]

% ==================================================
% REPOSITORIO
% ==================================================

{\fontsize{10.5}{13}\selectfont\bfseries
REPOSITORIO PÚBLICO
}\\[0.2cm]

\resizebox{0.92\textwidth}{!}{
\url{https://github.com/Melanie-G23/Prueba_Unidad_IV}
}

\vspace{0.5cm}

% ==================================================
% CAPTURA DEL RESUMEN
% ==================================================

{\fontsize{10.5}{13}\selectfont\bfseries
CAPTURA DEL RESUMEN DEL CUESTIONARIO SGA
}\\[0.15cm]

\includegraphics[
    width=0.72\textwidth,
    height=4cm,
    keepaspectratio
]{resumen_cuestionario.png}

\vspace{0.3cm}

% ==================================================
% CAPTURA DE EVALUACIÓN
% ==================================================

{\fontsize{10.5}{13}\selectfont\bfseries
CAPTURA DE LA EVALUACIÓN / INTENTO DEL SGA
}\\[0.15cm]

\includegraphics[
    width=0.72\textwidth,
    height=4cm,
    keepaspectratio
]{evaluacion_cuestionario.png}

\vfill

{\fontsize{10}{12}\selectfont
Quevedo -- Ecuador\\
2026
}

\end{titlepage}

% ==================================================
% P1. DIAGRAMA DE CLASES
% ==================================================

\section*{P1. Modelo de datos -- Diagrama de clases UML}

\begin{figure}[H]
    \centering
    \includegraphics[
        width=0.95\textwidth,
        keepaspectratio
    ]{diagrama_clases.png}
    \caption{Diagrama de clases UML del Sistema de Gestión de Pedidos.}
\end{figure}

% ==================================================
% P2. DIAGRAMA DE ACTIVIDADES
% ==================================================

\section*{P2. Modelo funcional -- Diagrama de actividades UML}

\begin{figure}[H]
    \centering
    \includegraphics[
        width=0.78\textwidth,
        keepaspectratio
    ]{diagrama_actividades.png}
    \caption{Diagrama de actividades del proceso Registrar pedido.}
\end{figure}

% ==================================================
% P3. DIAGRAMA DE ESTADOS
% ==================================================

\section*{P3. Modelo de comportamiento -- Máquina de estados UML}

\begin{figure}[H]
    \centering
    \includegraphics[
        width=\textwidth,
        keepaspectratio
    ]{diagrama_estado.png}
    \caption{Máquina de estados del ciclo de vida de un Pedido.}
\end{figure}

% ==================================================
% P4. CONSISTENCIA ENTRE PERSPECTIVAS
% ==================================================

\section*{P4. Consistencia entre perspectivas}

\begin{center}
\renewcommand{\arraystretch}{1.3}

\begin{tabular}{|p{3.5cm}|p{6cm}|p{3.5cm}|}
\hline
\textbf{Evento P3} & \textbf{Acción P2} & \textbf{Elemento P1} \\ \hline

registrarPedido &
Ingresar datos del pedido &
Pedido \\ \hline

verificarStock &
Verificar stock &
Producto.stock \\ \hline

confirmarPedido &
Confirmar pedido &
Pedido.estado \\ \hline

despacharPedido &
Despachar pedido &
Pedido.estado \\ \hline

entregarPedido &
Entregar pedido &
Pedido.estado \\ \hline

anularPedido &
Anular pedido &
Pedido.estado \\ \hline

\end{tabular}
\end{center}

\subsection*{Inconsistencia detectada}

En P3 existe el evento \texttt{anularPedido}, pero inicialmente esta
acción no se encontraba representada en el modelo funcional P2.

\subsection*{Corrección}

Se incorpora en P2 la acción ``Anular pedido'' con la condición
\texttt{[pedido no entregado]}, garantizando la consistencia entre
el modelo funcional y el modelo de comportamiento.

% ==================================================
% P5. REQUISITOS
% ==================================================

\section*{P5. Requisitos con esquema de atributos}

\subsection*{RF-01: Registrar pedido}

\textbf{Requisito:} El sistema deberá permitir al cliente registrar
un pedido compuesto por uno o más productos y sus respectivas cantidades.

\begin{itemize}
    \item \textbf{ID:} RF-01
    \item \textbf{Prioridad:} Alta
    \item \textbf{Estado:} Aprobado
    \item \textbf{Fuente:} Cliente
    \item \textbf{Estabilidad:} Alta
    \item \textbf{Versión:} 1.0
    \item \textbf{Verificación:} Prueba funcional
\end{itemize}

\subsection*{RF-02: Verificar stock}

\textbf{Requisito:} El sistema deberá verificar el stock de cada
producto antes de confirmar un pedido y, si no existe stock suficiente,
deberá notificar el faltante.

\begin{itemize}
    \item \textbf{ID:} RF-02
    \item \textbf{Prioridad:} Alta
    \item \textbf{Estado:} Aprobado
    \item \textbf{Fuente:} Tienda
    \item \textbf{Estabilidad:} Alta
    \item \textbf{Versión:} 1.0
    \item \textbf{Verificación:} Prueba funcional
\end{itemize}

\subsection*{RNF-01: Eficiencia de desempeño}

\textbf{Requisito:} El sistema deberá completar el registro y
confirmación de un pedido en menos de 2 segundos en al menos el
96\% de las operaciones.

\begin{itemize}
    \item \textbf{Característica ISO 25010:} Eficiencia de desempeño
    \item \textbf{Métrica:} Tiempo de respuesta
    \item \textbf{Umbral:} Menos de 2 segundos en al menos el 96\% de las operaciones
    \item \textbf{Prioridad:} Alta
    \item \textbf{Estado:} Aprobado
    \item \textbf{Fuente:} Tienda
    \item \textbf{Estabilidad:} Alta
    \item \textbf{Versión:} 1.0
    \item \textbf{Verificación:} Prueba de rendimiento
\end{itemize}

\subsection*{RNF-02: Seguridad}

\textbf{Requisito:} El sistema deberá impedir el acceso a los datos
personales del cliente a usuarios no autenticados en el 100\% de los
intentos de acceso evaluados.

\begin{itemize}
    \item \textbf{Característica ISO 25010:} Seguridad
    \item \textbf{Métrica:} Porcentaje de accesos no autorizados bloqueados
    \item \textbf{Umbral:} 100\%
    \item \textbf{Prioridad:} Alta
    \item \textbf{Estado:} Aprobado
    \item \textbf{Fuente:} Política de seguridad
    \item \textbf{Estabilidad:} Alta
    \item \textbf{Versión:} 1.0
    \item \textbf{Verificación:} Prueba de seguridad
\end{itemize}

% ==================================================
% P6. MOSCOW
% ==================================================

\section*{P6. Priorización MoSCoW}

\begin{center}
\renewcommand{\arraystretch}{1.3}

\begin{tabular}{|p{4cm}|p{2cm}|p{7cm}|}
\hline
\textbf{Requisito} & \textbf{MoSCoW} & \textbf{Justificación} \\ \hline

RF-01 Registrar pedido &
Must &
Es la función principal del sistema. \\ \hline

RF-02 Verificar stock &
Must &
Evita confirmar productos que no se encuentran disponibles. \\ \hline

RNF-01 Rendimiento &
Should &
Es importante para brindar una respuesta rápida. \\ \hline

RNF-02 Seguridad &
Must &
Es necesaria para proteger los datos de los usuarios. \\ \hline

\end{tabular}
\end{center}

% ==================================================
% P7. INSPECCIÓN 29148
% ==================================================

\section*{P7. Inspección 29148 (defectos -- retrabajo)}

\begin{center}
\small
\renewcommand{\arraystretch}{1.3}

\begin{tabular}{|c|c|c|c|c|c|p{3cm}|}
\hline
\textbf{Req.} &
\textbf{No ambiguo} &
\textbf{Completo} &
\textbf{Singular} &
\textbf{Verificable} &
\textbf{Factible} &
\textbf{Defecto} \\ \hline

RF-01 & Sí & Sí & Sí & Sí & Sí & Ninguno \\ \hline
RF-02 & No & Sí & No & Sí & Sí & Contiene verificar y notificar \\ \hline
RNF-01 & Sí & Sí & Sí & Sí & Sí & Ninguno \\ \hline
RNF-02 & Sí & Sí & Sí & Sí & Sí & Ninguno \\ \hline

\end{tabular}
\end{center}

\subsection*{Defecto}

RF-02 no cumple completamente la propiedad de singularidad porque
contiene dos obligaciones:

\begin{enumerate}
    \item Verificar stock.
    \item Notificar faltante.
\end{enumerate}

\subsection*{Retrabajo}

\textbf{RF-02:} El sistema deberá verificar la disponibilidad de stock
de cada producto antes de confirmar un pedido.

\medskip

\textbf{RF-03:} Cuando el stock disponible sea inferior a la cantidad
solicitada, el sistema deberá notificar al cliente el producto y la
cantidad faltante.

% ==================================================
% P8. PRUEBAS DE ACEPTACIÓN
% ==================================================

\section*{P8. Pruebas de aceptación trazadas}

\begin{center}
\small
\renewcommand{\arraystretch}{1.3}

\begin{tabular}{|c|c|p{2.2cm}|p{4cm}|p{5cm}|}
\hline
\textbf{Prueba} &
\textbf{Req.} &
\textbf{Tipo} &
\textbf{Entrada} &
\textbf{Criterio de éxito} \\ \hline

PA-01 &
RF-01 &
Funcional &
Cliente y productos &
Se registra correctamente el pedido con sus líneas. \\ \hline

PA-02 &
RF-02 &
Funcional &
Producto solicitado &
El sistema detecta stock insuficiente y no confirma el pedido. \\ \hline

PA-03 &
RNF-01 &
No funcional &
Registros de pedidos &
Al menos el 96\% se completa en menos de 2 segundos. \\ \hline

PA-04 &
RNF-02 &
No funcional &
Usuario no autenticado intenta acceder a datos &
El 100\% de accesos no autorizados son rechazados. \\ \hline

\end{tabular}
\end{center}

% ==================================================
% P9. MATRIZ DE TRAZABILIDAD
% ==================================================

\section*{P9. Matriz de trazabilidad}

\begin{center}
\small
\renewcommand{\arraystretch}{1.3}

\begin{tabular}{|c|p{3cm}|p{4.5cm}|c|c|}
\hline
\textbf{Requisito} &
\textbf{Fuente} &
\textbf{Modelo} &
\textbf{Prueba} &
\textbf{Dependencia} \\ \hline

RF-01 &
Cliente &
P1 Cliente/Pedido + P2 &
PA-01 &
RF-02 \\ \hline

RF-02 &
Tienda &
P1 Producto + P2 &
PA-02 &
RF-01 \\ \hline

RNF-01 &
Tienda &
Proceso Registrar pedido &
PA-03 &
RF-01 \\ \hline

RNF-02 &
Política de seguridad &
P1 Cliente &
PA-04 &
RF-01 \\ \hline

\end{tabular}
\end{center}

\subsection*{Trazabilidad}

\textbf{Hacia atrás:}

\[
\text{Requisito} \rightarrow \text{Fuente}
\]

\textbf{Hacia adelante:}

\[
\text{Requisito} \rightarrow \text{Modelo} \rightarrow \text{Prueba}
\]

\textbf{Horizontal:}

\[
\text{RF-01} \rightarrow \text{depende de RF-02}
\]

No se identifican requisitos huérfanos, debido a que los requisitos
poseen una fuente, un modelo asociado y una prueba de aceptación.

% ==================================================
% P10. GESTIÓN DEL CAMBIO
% ==================================================

\section*{P10. Gestión del cambio y línea base}

\subsection*{Solicitud de cambio RFC-01}

\textbf{Requisito afectado:} RNF-01

\textbf{Versión actual:} 1.0

\subsection*{Cambio solicitado}

Cambiar:

\begin{quote}
``menos de 2 segundos''
\end{quote}

por:

\begin{quote}
``menos de 1 segundo''.
\end{quote}

\textbf{Justificación:} Mejorar la rapidez percibida durante el
registro de pedidos.

\subsection*{Análisis de impacto}

Según la matriz de P9, el cambio afecta:

\begin{itemize}
    \item RNF-01.
    \item Proceso de registrar pedido.
    \item PA-03.
    \item Pruebas de rendimiento.
    \item Documentación de requisitos.
\end{itemize}

\subsection*{Decisión CCB}

\textbf{APROBADA.}

El CCB aprueba el cambio porque mejora la eficiencia de desempeño
y se considera técnicamente viable.

\subsection*{Nueva versión}

\[
\text{RNF-01 v1.0} \rightarrow \text{RNF-01 v1.1}
\]

\textbf{RNF-01 v1.1:} El sistema deberá completar el registro y
confirmación de un pedido en menos de 1 segundo en al menos el
96\% de las operaciones.

\subsection*{Línea base}

\textbf{Baseline v1.1}

\textbf{Configuración congelada:}

\begin{itemize}
    \item ERS v1.1
    \item Diagrama de clases v1.0
    \item Diagrama de actividades v1.0
    \item Máquina de estados v1.0
    \item Matriz de trazabilidad v1.1
    \item Pruebas de aceptación v1.1
    \item RFC-01 v1.0
\end{itemize}

\end{document}
